% Options for packages loaded elsewhere
\PassOptionsToPackage{unicode}{hyperref}
\PassOptionsToPackage{hyphens}{url}
\documentclass[
]{article}
\usepackage{xcolor}
\usepackage[margin=1in]{geometry}
\usepackage{amsmath,amssymb}
\setcounter{secnumdepth}{-\maxdimen} % remove section numbering
\usepackage{iftex}
\ifPDFTeX
  \usepackage[T1]{fontenc}
  \usepackage[utf8]{inputenc}
  \usepackage{textcomp} % provide euro and other symbols
\else % if luatex or xetex
  \usepackage{unicode-math} % this also loads fontspec
  \defaultfontfeatures{Scale=MatchLowercase}
  \defaultfontfeatures[\rmfamily]{Ligatures=TeX,Scale=1}
\fi
\usepackage{lmodern}
\ifPDFTeX\else
  % xetex/luatex font selection
\fi
% Use upquote if available, for straight quotes in verbatim environments
\IfFileExists{upquote.sty}{\usepackage{upquote}}{}
\IfFileExists{microtype.sty}{% use microtype if available
  \usepackage[]{microtype}
  \UseMicrotypeSet[protrusion]{basicmath} % disable protrusion for tt fonts
}{}
\makeatletter
\@ifundefined{KOMAClassName}{% if non-KOMA class
  \IfFileExists{parskip.sty}{%
    \usepackage{parskip}
  }{% else
    \setlength{\parindent}{0pt}
    \setlength{\parskip}{6pt plus 2pt minus 1pt}}
}{% if KOMA class
  \KOMAoptions{parskip=half}}
\makeatother
\setlength{\emergencystretch}{3em} % prevent overfull lines
\providecommand{\tightlist}{%
  \setlength{\itemsep}{0pt}\setlength{\parskip}{0pt}}
\usepackage{bookmark}
\IfFileExists{xurl.sty}{\usepackage{xurl}}{} % add URL line breaks if available
\urlstyle{same}
\hypersetup{
  hidelinks,
  pdfcreator={LaTeX via pandoc}}

\author{}
\date{}

\begin{document}

\section{La Historia de la Tesis Doctoral: Evaluación Cuantitativa de
Inteligencia Artificial Explicable
(XAI)}\label{la-historia-de-la-tesis-doctoral-evaluaciuxf3n-cuantitativa-de-inteligencia-artificial-explicable-xai}

\emph{Nivel Intermedio (Estudiantes Universitarios de Informática,
Ciencia de Datos e Ingeniería)}

\begin{center}\rule{0.5\linewidth}{0.5pt}\end{center}

\subsection{1. Introducción y Motivación: La Opacidad de la IA y el Reto
de la
Evaluación}\label{introducciuxf3n-y-motivaciuxf3n-la-opacidad-de-la-ia-y-el-reto-de-la-evaluaciuxf3n}

En la era del aprendizaje automático moderno, los modelos predictivos
como las Redes Neuronales Profundas (MLP), las Máquinas de Vectores de
Soporte (SVM) y los Ensambles de Árboles (Random Forest, XGBoost)
alcanzan un rendimiento excepcional. Sin embargo, su arquitectura
interna se comporta como una \textbf{``Caja Negra''}
(\(\mathbf{y} = f(\mathbf{x})\)), donde millones de parámetros hacen
imposible auditar de manera directa el proceso inferencial.

En sectores de alto impacto como las finanzas, la medicina o la
justicia, el marco legal y ético exige \textbf{explicabilidad} (Sistemas
de IA Explicable, XAI). Para lograrlo, la comunidad ha desarrollado
explicadores post-hoc agnósticos al modelo. Estos algoritmos intentan
aproximar el comportamiento de \(f(\mathbf{x})\) alrededor de una
instancia sin necesidad de acceder a la arquitectura interna del modelo.

Sin embargo, surgió un problema de investigación fundamental:
\textbf{¿Cómo sabemos si las explicaciones producidas por estos
explicadores son confiables?} Diferentes métodos de explicación (como
LIME o SHAP) aplicados al mismo modelo y a la misma instancia
frecuentemente producen explicaciones contradictorias. Sin un protocolo
estandarizado de evaluación cuantitativa, los científicos y auditores
corren el riesgo de caer en una ``ilusión de transparencia''.

Esta tesis aborda dicha brecha mediante el desarrollo, validación y
benchmarking del \textbf{Protocolo FOM-7} (\emph{Framework for
Operational Metrics in 7 Gates}).

\begin{center}\rule{0.5\linewidth}{0.5pt}\end{center}

\subsection{2. Los Métodos de Explicabilidad y las Métricas
Cuantitativas}\label{los-muxe9todos-de-explicabilidad-y-las-muxe9tricas-cuantitativas}

\subsubsection{A. Los Métodos de Explicabilidad
Estudiados}\label{a.-los-muxe9todos-de-explicabilidad-estudiados}

\begin{enumerate}
\def\labelenumi{\arabic{enumi}.}
\tightlist
\item
  \textbf{Atribución de Características (Feature Attribution):}

  \begin{itemize}
  \tightlist
  \item
    \textbf{LIME (Local Interpretable Model-agnostic Explanations):}
    Genera una perturbación sintética en el vecindario de la instancia
    \(\mathbf{x}\) y entrena un modelo sustituto lineal simple (\(g\))
    ponderado por proximidad espacial.
  \item
    \textbf{SHAP (Shapley Additive exPlanations):} Calcula la
    contribución marginal de cada característica basándose en la teoría
    de juegos cooperativos (Valores de Shapley), garantizando
    propiedades axiomáticas como consistencia y eficiencia.
  \end{itemize}
\item
  \textbf{Reglas Locales (Rule-based):}

  \begin{itemize}
  \tightlist
  \item
    \textbf{Anchors:} Construye reglas lógicas del tipo
    \texttt{SI\ condición\ ENTONCES\ predicción} con garantías
    estadísticas de precisión local en un vecindario de datos.
  \end{itemize}
\item
  \textbf{Explicaciones Contrafactuales (Counterfactuals):}

  \begin{itemize}
  \tightlist
  \item
    \textbf{DiCE (Diverse Counterfactual Explanations):} Encuentra las
    modificaciones mínimas necesarias en las características de la
    entrada para cambiar la predicción del modelo hacia una clase
    objetivo deseada.
  \end{itemize}
\end{enumerate}

\subsubsection{B. Las 5 Métricas Cuantitativas de
Evaluación}\label{b.-las-5-muxe9tricas-cuantitativas-de-evaluaciuxf3n}

Para comparar objetivamente estos métodos, la tesis formaliza cinco
métricas cuantitativas: 1. \textbf{Fidelidad (Fidelity):} Grado de
acuerdo entre la predicción del modelo sustituto/explicador y el modelo
caja negra original. 2. \textbf{Estabilidad (Stability):} Continuidad de
la explicación ante pequeñas perturbaciones en la entrada (medida
mediante distancias de repetición o coeficientes de variación). 3.
\textbf{Parsimonia (Sparsity / Complexity):} Fracción de características
no nulas utilizadas en la explicación (privilegiando explicaciones
breves y cognitivamente manejables). 4. \textbf{Coste Computacional
(Cost):} Tiempo de ejecución (\emph{wall-clock time}) en milisegundos
por instancia explicada. 5. \textbf{Brecha de Fidelidad (Faithfulness
Gap):} Error cuadrático medio entre la degradación real del modelo al
remover características y la expectativa teórica del explicador.

\begin{center}\rule{0.5\linewidth}{0.5pt}\end{center}

\subsection{3. La Estructura de la Investigación: Artículos A a
F}\label{la-estructura-de-la-investigaciuxf3n-artuxedculos-a-a-f}

La tesis doctoral se desarrolla a través de 6 artículos interconectados
que construyen progresivamente la evidencia científica:

\begin{verbatim}
[Opacidad de IA] ──> [Conflicto LIME vs SHAP] ──> [Necesidad de Gobernanza]
                                   │
  ┌────────────────────────────────┴────────────────────────────────┐
  │                                                                 │
  ▼                                                                 ▼
[Artículo A: Protocolo FOM-7] ──> [Artículo B: LIME vs SHAP] ──> [Artículo C: Escalabilidad]
  (Publicado - RIMI 2026)          (Enviado - CLEIej)              (Listo para envío)
                                                                    │
  ┌─────────────────────────────────────────────────────────────────┘
  │
  ▼
[Artículo D: Aplicación Empírica] ──> [Artículo E: Frontera Pareto] ──> [Artículo F: Alineación Humana]
  (Listo para envío - Tec. Marcha)      (Enviado - CyS CIC-IPN)          (En desarrollo - JCSI)
\end{verbatim}

\begin{center}\rule{0.5\linewidth}{0.5pt}\end{center}

\subsubsection{Artículo A: Fundamentación Teórica y el Protocolo
FOM-7}\label{artuxedculo-a-fundamentaciuxf3n-teuxf3rica-y-el-protocolo-fom-7}

\begin{itemize}
\tightlist
\item
  \textbf{Estado:} Publicado en \emph{Revista de Investigación
  Multidisciplinaria Iberoamericana (RIMI)} (2026, DOI:
  \texttt{10.69850/rimi.vi3.307}).
\item
  \textbf{Objetivo:} Definir un estándar metodológico auditable para
  evaluar explicadores post-hoc sin incurrir en fugas de datos o sesgos
  de varianza.
\item
  \textbf{Aporte:} Formalización de las \textbf{7 Puertas de FOM-7}:

  \begin{enumerate}
  \def\labelenumi{\arabic{enumi}.}
  \tightlist
  \item
    \emph{Congelación:} Pre-registro determinista del diseño
    experimental y del plan inferencial.
  \item
    \emph{Ejecución:} Generación de resultados crudos por celda
    experimental con semillas fijas.
  \item
    \emph{Auditoría:} Verificación de integridad de artefactos (filtrado
    de ejecuciones nulas o corruptas).
  \item
    \emph{Armonización:} Agregación de métricas por bloques
    homogeneizados (modelo \(\times\) tamaño muestra).
  \item
    \emph{Exportación:} Evaluación estadística no paramétrica (Pruebas
    de Friedman y Nemenyi con ajuste de Holm).
  \item
    \emph{Perfilado:} Aislamiento de la varianza algorítmica (semilla)
    frente a los efectos principales del método.
  \item
    \emph{Reporte:} Registro trazable de afirmaciones científicas
    delimitadas por constructo y dominio.
  \end{enumerate}
\end{itemize}

\begin{center}\rule{0.5\linewidth}{0.5pt}\end{center}

\subsubsection{Artículo B: El Duelo Comparativo --- LIME vs SHAP a
Escala}\label{artuxedculo-b-el-duelo-comparativo-lime-vs-shap-a-escala}

\begin{itemize}
\tightlist
\item
  \textbf{Estado:} Enviado a \emph{CLEI Electronic Journal (CLEIej)}
  (2026-10-04, Subm. \#1196).
\item
  \textbf{Objetivo:} Ejecutar una comparación empírica masiva entre los
  dos estándares de la industria (LIME y SHAP).
\item
  \textbf{Aporte:}

  \begin{itemize}
  \tightlist
  \item
    Demostración empírica de las ventajas teóricas de SHAP en
    estabilidad y fidelidad.
  \item
    Identificación de la fragilidad estocástica de LIME: al muestrear
    aleatoriamente perturbaciones locales, LIME exhibe alta variabilidad
    en los coeficientes asignados.
  \item
    LIME compensa esta limitación mediante un menor coste computacional
    y una mayor parsimonia en sus explicaciones.
  \end{itemize}
\end{itemize}

\begin{center}\rule{0.5\linewidth}{0.5pt}\end{center}

\subsubsection{Artículo C: Dinámica de Escalabilidad y
Resistencia}\label{artuxedculo-c-dinuxe1mica-de-escalabilidad-y-resistencia}

\begin{itemize}
\tightlist
\item
  \textbf{Estado:} Listo para envío a \emph{Tecnología en Marcha}
  (Edición Especial IA, Zenodo v0.12.0 DOI:
  \texttt{10.5281/zenodo.23165763}).
\item
  \textbf{Objetivo:} Analizar cómo se comportan los explicadores cuando
  varía el tamaño de muestra del conjunto de datos y la profundidad de
  los modelos.
\item
  \textbf{Aporte:} Cuantificación del punto de degradación de los
  sustitutos locales. Se demuestra que incrementar el tamaño del
  conjunto de datos de fondo mejora la estabilidad de KernelSHAP a costa
  de un crecimiento cuadrático en el tiempo de procesamiento.
\end{itemize}

\begin{center}\rule{0.5\linewidth}{0.5pt}\end{center}

\subsubsection{Artículo D: Benchmark Empírico Aplicado (UCI Adult
Income)}\label{artuxedculo-d-benchmark-empuxedrico-aplicado-uci-adult-income}

\begin{itemize}
\tightlist
\item
  \textbf{Estado:} Listo para envío a \emph{Tecnología en Marcha}
  (Edición Especial IA).
\item
  \textbf{Objetivo:} Aplicar el protocolo FOM-7 sobre un caso de estudio
  de alta relevancia social y económica (clasificación de ingresos
  tabulares \emph{UCI Adult Income}).
\item
  \textbf{Aporte:}

  \begin{itemize}
  \tightlist
  \item
    Evaluación sobre 300 celdas planificadas
    (\(5 \text{ modelos} \times 4 \text{ métodos} \times 5 \text{ semillas} \times 3 \text{ tamaños}\)),
    obteniendo 275 celdas calificadas tras la auditoría FOM-7.
  \item
    Comparación de 5 familias de modelos predictivos: Logistic
    Regression, Random Forest, XGBoost, SVM (RBF) y MLP.
  \item
    Confirmación de que las explicaciones estructurales (Anchors) y
    contrafactuales (DiCE) deben evaluarse en dimensiones distintas a la
    simple atribución de características.
  \end{itemize}
\end{itemize}

\begin{center}\rule{0.5\linewidth}{0.5pt}\end{center}

\subsubsection{Artículo E: Fronteras de Pareto y Gobernanza de
XAI}\label{artuxedculo-e-fronteras-de-pareto-y-gobernanza-de-xai}

\begin{itemize}
\tightlist
\item
  \textbf{Estado:} Enviado a \emph{Computación y Sistemas (CIC-IPN)}
  (2026-10-04, Subm. \#6783).
\item
  \textbf{Objetivo:} Resolver el problema de selección de explicadores
  mediante optimización multiobjetivo y análisis de frontera de
  eficiencia (Pareto).
\item
  \textbf{Aporte:}

  \begin{itemize}
  \tightlist
  \item
    Demostración de que ningún método domina simultáneamente en todas
    las métricas.
  \item
    Formulación de criterios de gobernanza:

    \begin{itemize}
    \tightlist
    \item
      \textbf{Entornos de Alta Severidad (Auditoría Médica/Legal):}
      Selección obligatoria de \textbf{SHAP} debido a su consistencia
      axiomática y alta estabilidad.
    \item
      \textbf{Entornos de Tiempo Real / Producción de Alta Velocidad:}
      Selección recomendada de \textbf{LIME} o \textbf{TreeSHAP} para
      equilibrar coste y latencia.
    \end{itemize}
  \end{itemize}
\end{itemize}

\begin{center}\rule{0.5\linewidth}{0.5pt}\end{center}

\subsubsection{Artículo F: Alineación Semántica y Evaluación por
Humanos/LLM}\label{artuxedculo-f-alineaciuxf3n-semuxe1ntica-y-evaluaciuxf3n-por-humanosllm}

\begin{itemize}
\tightlist
\item
  \textbf{Estado:} En desarrollo inicial (\emph{Journal of Computer
  Sciences Institute}).
\item
  \textbf{Objetivo:} Conectar las métricas computacionales objetivas con
  la percepción humana de explicabilidad y la evaluación mediante
  Modelos de Lenguaje (LLM-as-a-Judge).
\item
  \textbf{Aporte:} Protocolo de validación cruzada para verificar si los
  rankings matemáticos de fidelidad y estabilidad se traducen en una
  mejor comprensión y confianza por parte de los usuarios finales.
\end{itemize}

\begin{center}\rule{0.5\linewidth}{0.5pt}\end{center}

\subsection{4. Síntesis y Relevancia para la Defensa de
Tesis}\label{suxedntesis-y-relevancia-para-la-defensa-de-tesis}

La narrativa unificada de esta tesis doctoral demuestra que la
evaluación de XAI no es una búsqueda de ``el mejor método único'', sino
la construcción de un \textbf{sistema de perfiles explicativos
trazables}:

\begin{enumerate}
\def\labelenumi{\arabic{enumi}.}
\tightlist
\item
  \textbf{Problema de Investigación:} La falta de rigor y la
  contradicción entre explicadores agnósticos.
\item
  \textbf{Solución Metodológica:} El protocolo \textbf{FOM-7} como marco
  de gobernanza y control de varianza.
\item
  \textbf{Evidencia Empírica (Papeles B, C, D):} Cuantificación precisa
  del comportamiento de LIME, SHAP, Anchors y DiCE.
\item
  \textbf{Aplicación Práctica (Papeles E, F):} Guía de selección basada
  en fronteras de Pareto y alineación con las necesidades del usuario.
\end{enumerate}

Este trabajo transforma la explicabilidad en una evidencia científica
reproducible, defendible y auditable.

\end{document}
